\documentclass[11pt]{article}
\usepackage[T1]{fontenc}
\usepackage{lmodern,amsmath,amssymb,amsthm,booktabs,microtype}
\usepackage[a4paper,margin=25mm]{geometry}
\usepackage{xurl}
\usepackage[hidelinks]{hyperref}
\newtheorem{theorem}{Theorem}[section]
\newtheorem{lemma}[theorem]{Lemma}
\newtheorem{proposition}[theorem]{Proposition}
\newtheorem{corollary}[theorem]{Corollary}
\newcommand{\Q}{\mathbb Q}
\newcommand{\Z}{\mathbb Z}
\newcommand{\C}{\mathbb C}
\newcommand{\OO}{\mathcal O}
\newcommand{\N}{\operatorname N}
\newcommand{\cont}{\operatorname{cont}}
\newcommand{\workbenchurl}{https://github.com/KokunoYumeto/zeta-and-l-functions-workbench}
\newcommand{\recordurl}{https://github.com/KokunoYumeto/mathematics-workbenches}
\title{Rational content and the cubic Riemann hypothesis, after Hecke}
\author{}
\date{24 November 2026}
\begin{document}
\maketitle
\begin{abstract}
Hecke's theory of ideal characters supplies the cubic $L$-function
studied here. The decomposition of a quadratic lattice vector into
its rational content and a primitive vector separates one tractable
part of its prime-extraction problem. We apply it to the
faithful cubic ray character of conductor $(269)$ over the
Eisenstein integers. A smooth restriction to rational content
greater than $R$ has the bound
$O_\epsilon(Y^{1+\epsilon}R^{-3/2})$ in a norm window of size $Y$.
Classical RH for the Riemann zeta function improves the exponent
of $R$ to $-7/4$. Conditional on a displayed logarithmic
convolution identity, this permits deletion of contents at least
$Y^{1/3}$ without a zero-location hypothesis, or $Y^{2/7}$
under zeta RH, at the
square-root scale. The retained primitive range remains long.
We give its exact arithmetic coefficients and show why neither
periodicity nor the current cubic large sieve bounds it at that
scale. Exact rational-prime and lattice computations agree
through norm $200000$. No cubic RH estimate is proved.
The mathematics and text were produced by GPT-6 Astra in ChatGPT
Pro mode; no human has checked this text.
\end{abstract}

\section{The content estimate}

Put $K=\Q(\omega)$, $\omega^2+\omega+1=0$, and
$\OO=\Z[\omega]$. Delete the prime ideals above $3$ and $269$;
all ideal sums below use only the remaining ideals.
Let $\psi$ be either faithful cubic character of
$K(\sqrt[3]{269})/K$. Its ray-class description has conductor
$(269)$ and gives, up to conjugation,
\[
 \psi((a+b\omega))=\xi(a+b\omega\bmod269),\qquad
 \xi(v)\ \longleftrightarrow\ v^{24120}\in\mu_3
       \subset\mathbb F_{269^2}.
                                                        \tag{1.1}
\]
In particular $\psi$ is trivial on rational integers prime to
$269$. Indeed $24120=90(269-1)$, and it kills the six Eisenstein
units as well. The ray-class identification uses the class-number-one
ray sequence and class field correspondence
(Milne, 2020, V.1.7 and V.3.5--3.7).
These classical identifications specify the fixed character used below.

For a nonzero ideal $\mathfrak a=(a+b\omega)$ define
\[
                 \cont(\mathfrak a)=\gcd(a,b)>0.
\]
Multiplication by a unit acts through an invertible integer matrix
on $(a,b)$, so the definition is independent of the generator.
Write $\mu_K$ for ideal M\"obius, $\mu$ for integer M\"obius,
and $\eta=\chi_{-3}$. Let $\chi:[0,\infty)\to[0,1]$ be smooth,
equal to one on $[0,1]$ and zero on $[2,\infty)$.
For $V\in C_c^\infty((1,2))$, $Y\ge2$, and $R\ge1$, define
\[
 M_{V,\chi}(Y,R)=
 \sum_{\mathfrak a}\mu_K(\mathfrak a)\psi(\mathfrak a)
 V(\N\mathfrak a/Y)
 \bigl(1-\chi(\cont(\mathfrak a)/R)\bigr).                 \tag{1.2}
\]

\begin{theorem}\label{thm:content}
For every $\epsilon>0$,
\[
 M_{V,\chi}(Y,R)
       \ll_{V,\chi,\epsilon}Y^{1+\epsilon}R^{-3/2}.        \tag{1.3}
\]
If classical RH for $\zeta(s)$ is granted, then
\[
 M_{V,\chi}(Y,R)
       \ll_{V,\chi,\epsilon}Y^{1+\epsilon}R^{-7/4}.        \tag{1.4}
\]
The constants are uniform over bounded families of smooth $V$
with the same fixed compact support in $(1,2)$; for each
$\epsilon$ only finitely many smooth seminorms are needed.
The second assertion is a smooth-cutoff statement, not an
unsmoothed content-tail estimate.
\end{theorem}

The absolute estimate for the same range is
$O(Y/R)$, from lattice counting. Thus (1.3) saves a factor
$R^{1/2}$ on that particular range, and (1.4) saves $R^{3/4}$.
This is a sparse portion of the complete inverse sum, not a bound
for all of it. The argument uses that the particular cubic
character is trivial on rational content.

\subsection{Application to the logarithmic convolution}

The following application is conditional on the displayed
convolution identity. Suppose that, for a fixed
$W\in C_c^\infty((1,2))$, fixed $\kappa>2$, and
\[
 L=(\log(3X))^2(\log(e+\log(3X)))^\kappa,
\]
the following holds for every $M>0$:
\[
 T_W(X)=R_W(X;L)+O_M(X^{-M}),                             \tag{4.1}
\]
where
\begin{align}
 T_W(X)&=\sum_{\mathfrak n}\Lambda_K(\mathfrak n)
               \psi(\mathfrak n)W(\N\mathfrak n/X),\notag\\
 R_W(X;L)&=\sum_{\N\mathfrak c<2L}\psi(\mathfrak c)
       \log\N\mathfrak c
       \sum_{\mathfrak d}\mu_K(\mathfrak d)\psi(\mathfrak d)
       W(\N\mathfrak c\N\mathfrak d/X)
       (1-\chi(L\N\mathfrak d/X)).                       \tag{4.2}
\end{align}
The inverse variable satisfies $\N\mathfrak d>X/L$ when
the cutoff is nonzero.

Put $Y_{\mathfrak c}=X/\N\mathfrak c$. Define
$R_{\rm low}(X;L,\gamma)$ by multiplying the inner summand
in (4.2) by
\[
                 \chi\bigl(\cont(\mathfrak d)/
                                  Y_{\mathfrak c}^{\gamma}\bigr).
\]
This is an exact smooth partition.

\begin{corollary}\label{cor:remainder}
Assume the identity (4.1). Then for every $\epsilon>0$,
\[
 T_W(X)=R_{\rm low}(X;L,1/3)+O_\epsilon(X^{1/2+\epsilon})
                                                        \tag{4.3}
\]
without a zero-location hypothesis. Under classical zeta RH,
\[
 T_W(X)=R_{\rm low}(X;L,2/7)+O_\epsilon(X^{1/2+\epsilon}).
                                                        \tag{4.4}
\]
Neither assertion estimates $R_{\rm low}$.
\end{corollary}

\section*{Proof overview}
For a primitive Eisenstein vector of norm $q$, the ideal M\"obius
value is the ordinary $\mu(q)$. Rational content contributes
$\mu(r)^2\eta(r)$, with the essential condition $(r,q)=1$.
We first sum in $r$. The squarefree expansion and the bounded
partial sums of $\eta$ give square-root cancellation in its
length. Under zeta RH, its rational Dirichlet series gives a
quarter-power smooth estimate. Counting the remaining primitive
vectors proves (1.3)--(1.4). We keep the moving content cutoff
inside the smooth weight throughout.

\section{Exact arithmetic and the elementary bound}

\begin{lemma}\label{lem:identity}
Write $\alpha=r\alpha_0$, where $r=\cont((\alpha))$ and
$\alpha_0=a+b\omega$ has $\gcd(a,b)=1$. Put
$q=\N\alpha_0=a^2-ab+b^2$. Away from the deleted primes,
\[
 \mu_K((\alpha))\psi((\alpha))
   =\mu(r)^2\eta(r)\mu(q)\psi((\alpha_0))
                      {\bf1}_{(r,q)=1}.                 \tag{2.1}
\]
This identity includes all zero cases.
\end{lemma}

\begin{proof}
If a rational prime $p\equiv2\pmod3$ divides $q$, it divides
both coordinates of $\alpha_0$, since $p$ is inert.
This contradicts primitivity. If $p\equiv1\pmod3$ splits as
$\mathfrak p\overline{\mathfrak p}$, both factors cannot divide
$(\alpha_0)$, because their product is $(p)$.
Consequently the valuation of $q$ at a split prime equals the
valuation of its unique possible prime-ideal factor in
$(\alpha_0)$. The prime above three has been excluded.
It follows that
\[
                   \mu_K((\alpha_0))=\mu(q).
\]
In particular a nonzero value forces $q$ squarefree and every
prime factor of $q$ to be $1\pmod3$.

For the rational content, a squarefree $r$ contributes two prime
ideals at each split prime and one at each inert prime. Therefore
$\mu_K((r))=\eta(r)$ when $r$ is squarefree. A squared rational
prime factor makes this ideal M\"obius value zero.
A prime common to $r$ and $q$ likewise creates a repeated
prime-ideal factor. These are precisely the additional zero
conditions. Finally $\psi((r))=1$ by (1.1). Multiplicativity
proves (2.1).
\end{proof}

The lattice interpretation makes every factor visible:
\begin{align}
 M_{V,\chi}(Y,R)
 =\frac16\sum_{\substack{\alpha_0=a+b\omega\ne0\\
                         (a,b)=1\\(q,3\cdot269)=1}}
 &\mu(q)\psi((\alpha_0))
 \sum_{\substack{r\ge1\\(r,269q)=1}}
      \mu(r)^2\eta(r) \notag\\[-2pt]
 &{}\times V(qr^2/Y)(1-\chi(r/R)).                        \tag{2.2}
\end{align}
The factor $1/6$ accounts for the six associates.
The factor $\eta$ already vanishes at multiples of three.

\begin{lemma}\label{lem:scalar}
For every positive integer $h$ prime to three,
\[
 \sum_{\substack{r\le T\\(r,h)=1}}\mu(r)^2\eta(r)
                      \ll \sqrt T\,\tau(h),\qquad T\ge1. \tag{2.3}
\]
Consequently, for $F$ in a bounded family of smooth functions
supported in a fixed compact interval in $(0,\infty)$,
\[
 \sum_{\substack{r\ge1\\(r,h)=1}}\mu(r)^2\eta(r)F(r/T)
                      \ll_F \sqrt T\,\tau(h).            \tag{2.4}
\]
\end{lemma}

\begin{proof}
Let $A(t)=\sum_{1\le n\le t}\eta(n)$; then $|A(t)|\le1$.
Insert $\mu(r)^2=\sum_{d^2\mid r}\mu(d)$, write $r=d^2u$,
and expand the remaining coprimality by its squarefree divisors.
The exact result is
\[
 \sum_{\substack{d\le\sqrt T\\(d,3h)=1}}\mu(d)
     \sum_{e\mid\operatorname{rad}(h)}
          \mu(e)\eta(e) A\left(\frac{T}{d^2e}\right).
\]
It has absolute value at most $\lfloor\sqrt T\rfloor
\tau(\operatorname{rad}h)\le\sqrt T\,\tau(h)$.
Partial summation gives (2.4), with a constant controlled by
the supremum and total variation of $F$.
\end{proof}

For each $\alpha_0$ in (2.2), put
\[
 T_q=\sqrt{Y/q},\qquad
 F_{q,Y,R}(v)=V(v^2)\bigl(1-\chi((T_q/R)v)\bigr).
                                                        \tag{2.5}
\]
These weights have fixed support in $1<v<\sqrt2$ and uniformly
bounded derivatives of every fixed order.
Indeed, when $T_q/R\ge2$ the second factor is identically one
on the support; when $T_q/R\le1/\sqrt2$ it is identically zero.
All nonconstant cases have $T_q/R$ in the compact interval
$[1/\sqrt2,2]$. This observation also shows that a nonzero
inner sum requires
\[
                       q<2Y/R^2.                       \tag{2.6}
\]
Values $T_q\in[1/\sqrt2,1]$ can be included in (2.4) by changing
its absolute constant.

The number of Eisenstein vectors of norm at most $Z$ is $O(Z)$
for $Z\ge1$. For example, the norm is bounded below by
$(a^2+b^2)/2$, which gives this estimate by an enclosing square.
Together with $\tau(269q)\ll_\epsilon q^\epsilon$,
(2.2)--(2.6) imply
\begin{align*}
 |M_{V,\chi}(Y,R)|
 &\ll_\epsilon
       Y^{1/4}\sum_{\N\alpha_0<2Y/R^2}
                    (\N\alpha_0)^{-1/4+\epsilon}\\
 &\ll_\epsilon
       Y^{1/4}(Y/R^2)^{3/4+\epsilon}\\
 &\ll_\epsilon Y^{1+\epsilon}R^{-3/2}.
\end{align*}
If $2Y/R^2<1$, the sum is empty. This proves (1.3).

\section{The zeta-RH refinement, with all local factors}

Only in this section is classical RH for $\zeta$ assumed.
For $(h,3)=1$, the rational Dirichlet series of the scalar
sequence in Lemma~\ref{lem:scalar} is
\[
 F_h(s)=\sum_{\substack{r\ge1\\(r,h)=1}}
             \frac{\mu(r)^2\eta(r)}{r^s}
 =\frac{L(s,\eta)}
        {\zeta(2s)(1-3^{-2s})}
      \prod_{p\mid h}(1+\eta(p)p^{-s})^{-1},
          \qquad\Re s>1.                               \tag{3.1}
\]
The factor at three is essential: $\eta(3)=0$, so its squarefree
Euler factor is one. The restriction at 269, and at every
prime dividing the primitive norm, occurs in the final product.

For every fixed $\sigma_0>0$ and $\epsilon>0$,
\[
 \sup_{\Re s\ge\sigma_0}
 \prod_{p\mid h}|1+\eta(p)p^{-s}|^{-1}
                      \ll_{\sigma_0,\epsilon} h^\epsilon. \tag{3.2}
\]
To prove this, bound each inverse by $(1-p^{-\sigma_0})^{-1}$.
For all sufficiently large primes, its logarithm is at most
$\epsilon\log p$. The remaining finitely many primes contribute
a constant depending only on $\sigma_0,\epsilon$.
Thus the estimate is uniform in the potentially growing
integer $h=269q$, not just for fixed punctures.

Under zeta RH the right side of (3.1) is holomorphic on
$\Re s>1/4$. The pole of $\zeta$ at $2s=1$ gives a zero of its
reciprocal and causes no residue. The local factors have their
possible zeros or poles only on $\Re s=0$.
On every fixed closed inner strip, $1/\zeta(2s)$ has polynomial
growth in $|\Im s|$. One short justification is to remove the
pole from $\zeta(z)$ by using $(z-1)\zeta(z)$, take its analytic
logarithm in the zero-free half-plane $\Re z>1/2$, and apply
Borel--Carath\'eodory on disks centered at $2+it$.
The Euler product anchors the logarithm there; the functional
equation and ordinary strip bounds give an upper real part
$O(\log(3+|t|))$ on a slightly larger inner strip.
The reciprocal therefore has polynomial size on a smaller one.
This argument grants zeta RH; it does not import cubic RH.

Mellin inversion and a shift to $\Re s=1/4+\delta$, for fixed
$\delta>0$, now give the uniform smooth estimate
\[
 \sum_{\substack{r\ge1\\(r,h)=1}}\mu(r)^2\eta(r)F(r/T)
       \ll_{\delta,\epsilon,F}
                        T^{1/4+\delta}h^\epsilon.        \tag{3.3}
\]
Rapid Mellin decay dominates the stated polynomial growth.
There are no crossed poles. Uniformity holds for bounded
smooth families, hence for exactly the weights (2.5).
No characteristic function of a sharp interval is substituted
for those weights.

Apply (3.3) inside (2.2), choosing $0<\delta<1/4$ and
putting $\theta=1/4+\delta$.
Lattice counting yields
\[
 |M_{V,\chi}(Y,R)|
 \ll_{\delta,\epsilon}
     Y^{\theta/2}(Y/R^2)^{1-\theta/2+\epsilon}
 =Y^{1+\epsilon}R^{-7/4+\delta-2\epsilon}.                \tag{3.4}
\]
Whenever the sum is nonzero, $R\le\sqrt{2Y}$.
Choose the internal positive $\delta,\epsilon$ sufficiently
small in terms of the desired final error exponent, and
absorb $R^\delta$ into $Y^{\delta/2}$.
This proves (1.4).

\section{Effect on the logarithmic convolution}



\begin{proof}[Proof of Corollary~\ref{cor:remainder}]
For each direct ideal $\mathfrak c$, the inner smooth weight is
\[
 V_{\mathfrak c}(u)=W(u)
            (1-\chi((L/\N\mathfrak c)u)).
\]
As $\mathfrak c$ varies, these weights have uniformly bounded
smooth seminorms: for $L/\N\mathfrak c\ge2$ the cutoff is
identically zero in the subtracted term, and the remaining
parameter values lie in $[1/2,2]$.
Apply Theorem~\ref{thm:content} at
$R=Y_{\mathfrak c}^{1/3}$, or at
$R=Y_{\mathfrak c}^{2/7}$ in its conditional part.
Each high-content inner sum is
$O_\epsilon(Y_{\mathfrak c}^{1/2+\epsilon})$.
Counting direct ideals gives
\[
 |R_W-R_{\rm low}|
 \ll_\epsilon X^{1/2+\epsilon}
       \sum_{\N\mathfrak c<2L}
         \frac{\log\N\mathfrak c}
              {(\N\mathfrak c)^{1/2+\epsilon}}
 \ll_\epsilon X^{1/2+\epsilon}L^{1/2}\log(2L).
\]
Choose the internal error exponent smaller than the stated one.
Since $L$ is polylogarithmic, this proves the result using
(4.1). For sufficiently large $X$, all relevant
$Y_{\mathfrak c}\ge2$; the initial bounded range is absorbed
by the constant.
\end{proof}

Write $\mathfrak d=(r\alpha_0)$ in the retained range.
The exact remaining factor ranges are
\[
 \N\mathfrak c<2L,\quad
 Y_{\mathfrak c}<\N\mathfrak d<2Y_{\mathfrak c},\quad
 r<2Y_{\mathfrak c}^{\gamma},\quad
 \frac14Y_{\mathfrak c}^{1-2\gamma}
                 <\N\alpha_0<2Y_{\mathfrak c}.           \tag{4.5}
\]
Thus the primitive norm is at least
$Y_{\mathfrak c}^{1/3}/4$ in (4.3), or
$Y_{\mathfrak c}^{3/7}/4$ in (4.4).
These ranges are still polynomially long. In particular
$r=1$ has not been removed.

\section{The retained arithmetic and the methods still lacking a bound}

\subsection{The primitive coefficient and a fixed-period counterexample}

Let
\[
 h(n)=\sum_{\substack{\N\mathfrak a=n\\\cont(\mathfrak a)=1}}
                         \mu_K(\mathfrak a)\psi(\mathfrak a).
\]
It is multiplicative and supported on squarefree integers
whose prime factors are $1\pmod3$, away from the deleted primes.
At a split prime,
\[
 h(p)=-t(p),\qquad
 t(p)=\#\{z\bmod p:z^3=269\}-1\in\{2,-1\}.               \tag{5.1}
\]
Indeed the two primitive prime ideals contribute
$-\psi(\mathfrak p)-\psi(\overline{\mathfrak p})$.
At an inert prime there is no primitive ideal of its norm.
Equivalently, the retained lattice weight is
\[
       \mu(a^2-ab+b^2)\,\xi(a+b\omega\bmod269).            \tag{5.2}
\]
The ordinary M\"obius factor is not periodic with the
fixed ray character.

There is a small exact witness even after adjoining the local
period at three. The vectors $1+\omega$ and $1615+\omega$
are congruent modulo $807\OO$ and have the same character
value one. They are primitive and have norms
\[
                    1,\qquad
       2606611=7\cdot73\cdot5101,
\]
respectively. All three displayed factors are prime and
$1\pmod3$. Their ideal M\"obius values are therefore
$+1$ and $-1$. A complete period of the character cannot
be discarded after replacing its coefficient by (5.2).

For another precise check on the retained difficulty, put
$P(s)=L_K^S(s,\psi)$. Initially on $\Re s>1$,
\[
 H(s):=\sum_n h(n)n^{-s}=\frac{C(s)}{P(s)},                \tag{5.3}
\]
where
\[
 C(s)=
 \prod_{\substack{p\equiv2\ (3)\\p\ne269}}
       (1-p^{-2s})^{-1}
 \prod_{p\equiv1\ (3)}
       \frac{1-t(p)p^{-s}}
            {1-t(p)p^{-s}+p^{-2s}}.                     \tag{5.4}
\]
This Euler product converges absolutely and is nowhere zero
on $\Re s>1/2$. The factors differ from one by
$O(p^{-2\Re s})$ on compact inner strips.
The denominators at split primes factor as
$(1-zp^{-s})(1-\bar z p^{-s})$, $|z|=1$.
For a numerator with $t(p)=2$, nonvanishing follows from
$2p^{-\Re s}<1$, since every such prime is at least seven.
The other numerator is $1+p^{-s}$.
Thus every possible zero pole of $P$ in $\Re s>1/2$
survives already in this primitive coefficient series.

Equation (5.3) is an obstruction check, not a claimed new
method of estimating it. The primitive sum is a literal
piece of the restricted remainder: for $\mathfrak c=(2)$, which
has norm four and character one, and sufficiently large $X$,
the $r=1$ part of its summand is exactly
\[
                       (\log4)\sum_n h(n)W(4n/X).
                                                        \tag{5.5}
\]
Other direct ideals could cancel it, so this observation is
not a lower bound for the complete remainder.
It specifies a still-unestimated arithmetic component.

\subsection{What lattice and bilinear literature actually supplies}

Lam, Schindler, and Xiao (2020, Lemma~5.2), drawing on
Helfgott's quadratic-form argument, give
$O_A(Y(\log Y)^{-A})$ for sums of $\mu(Q(a,b))$ in fixed
residue classes, sectors, and polylogarithmic congruence
moduli, for a primitive positive definite quadratic form.
Finite residue expansion and partial summation apply this
to the fixed-period primitive weight (5.2).
This is genuine cancellation but does not reach the
square-root scale. Their modulus range also does not
automatically cover a coprimality modulus as large as the
remaining $r$ in (4.5).
This is an external estimate; the content bounds above do not depend on it.

For comparison with a cubic large-sieve estimate,
de Faveri (2026, Theorem~1.1) proves the fixed-order large sieve
for power-free rows and columns. At order three its bound is
\[
 \sum_{\N\mathfrak a\le A}^{\rm sf}
 \left|\sum_{\N\mathfrak b\le B}^{\rm sf}
       u_{\mathfrak b}\chi_{\mathfrak a}(\mathfrak b)\right|^2
 \ll_\epsilon(AB)^\epsilon
       \bigl(A+B+(AB)^{2/3}\bigr)\sum_{\mathfrak b}|u_{\mathfrak b}|^2.
                                                        \tag{5.6}
\]
It does not remove the cubic term already present for squarefree
indices. For the single row of norm $269^2$, fixed $A$, and
a column of length $B$, (5.6) gives only
$O_\epsilon(B^{1+\epsilon})$ for an unnormalized sum with
bounded coefficients. The target is $B^{1/2+\epsilon}$.
This large-sieve theorem is used only for the stated comparison of scales.

In (4.2), the short direct ideal does not supply the variable
row of (5.6). Multiplicativity gives
$\psi(\mathfrak c\mathfrak d)=\psi(\mathfrak c)\psi(\mathfrak d)$.
Cubic reciprocity replaces a symbol against the fixed element
269 by another symbol against that same fixed element; it
does not create $\chi_{\mathfrak c}(\mathfrak d)$.
On a rectangle of ideals, choosing bounded coefficient arrays
$u_{\mathfrak c}=\overline{\psi(\mathfrak c)}$ and
$v_{\mathfrak d}=\overline{\psi(\mathfrak d)}$ makes the
corresponding bilinear sum equal the positive size of that rectangle.
The content deletion leaves this issue intact.
A successful estimate must exploit the particular
M\"obius-weighted primitive sum or prove additional
oscillation in a valid transformation.

\section{Exact verification and scope}

The original program \emph{Cubic rational content certificate}
computes reciprocal norm coefficients in two ways.
The first uses rational-prime Euler factors and the test
$269^{(p-1)/3}\bmod p$ at split primes. The second enumerates
Eisenstein vectors and evaluates (2.1) through the finite field
$\mathbb F_{269^2}$. It sums the trace of the character and
divides by twelve, accounting for the six associates and
conjugation. It also retains the rational content in both
computations.

Through norm $200000$, all $200000$ norm coefficients agree,
as do all $23157$ nonzero content-refined entries.
The enumeration uses $483726$ lattice vectors away from
the excluded primes; $423324$ have nonzero ideal M\"obius.
Eighty-eight further exact checks expand the scalar sum
in Lemma~\ref{lem:scalar} and check its elementary bound.
These calculations verify conventions and finite identities.
They are not evidence sufficient to prove an asymptotic
square-root estimate.

The content estimate removes, at square-root cost, the smooth
large-content portion of the explicit long inverse remainder.
The unhandled
part is specified in (4.5), includes the primitive component
(5.5), and has no square-root bound here. Zeta RH is used
only in the explicitly conditional refinement.
No general host construction, abelian reduction of the cubic
function, or cubic RH theorem is asserted.

\section*{Method and availability}
The mathematics, text, and original program were produced by GPT-6 Astra
in ChatGPT Pro mode. The character is specified by the cited ray-class
description. The content identity, uniform scalar estimates, and
conditional convolution consequence are proved here. Classical zeta RH
is used only for the explicitly conditional refinement. The cited
quadratic-form and large-sieve estimates serve as comparisons, not as
proofs of a retained square-root bound. The accompanying records distinguish
finite computation from targeted or statement-only mathematical checks.
No human review or formal verification is claimed.
The original text and program are dedicated to the public domain under
CC0 1.0; cited works retain their own terms. Sources and verification
records are at \url{\workbenchurl}, and version records are at
\url{\recordurl}.

\begin{thebibliography}{9}
\small
\bibitem{Milne}
Milne, J. S. (2020). \emph{Class field theory} (Version 4.03).
V.1.7, p.~150; V.3.5--3.7, p.~158.
\url{https://www.jmilne.org/math/CourseNotes/CFT.pdf}.

\bibitem{Poonen}
Poonen, B. (2015). \emph{Tate's thesis}.
MIT 18.786 lecture notes, Theorems 5.16 and 5.22, pp.~34--39.
\url{https://math.mit.edu/~poonen/786/notes.pdf}.

\bibitem{LSX}
Lam, P. C.-H., Schindler, D., \& Xiao, S. Y. (2020).
On prime values of binary quadratic forms with a thin variable.
\emph{Journal of the London Mathematical Society, 102}(2),
749--772. Lemma 5.2 and its preceding discussion; the same
statement is in arXiv:1809.10755v1, pp.~25--26.
\url{https://doi.org/10.1112/jlms.12336}.

\bibitem{dF}
de Faveri, A. (2026, October 2).
\emph{Optimal large sieve for fixed order characters}
(arXiv:2610.04045v1). Theorem 1.1, p.~2.
\url{https://arxiv.org/abs/2610.04045v1}.
\end{thebibliography}
\end{document}

